\chapter{Kampfregeln}

Dieses Kapitel beschreibt alle Regeln, die im Kampf relevant sind. Kämpfe laufen
rundenbasiert ab und bestehen aus einer festen Abfolge von Schritten.

\section{Der Ablauf eines Kampfes}

\begin{rulebox}{Kampf in Kurzform}
	Ein Kampf besteht aus aufeinanderfolgenden \textbf{Runden}.  
	In jeder Runde handeln alle Beteiligten genau einmal in festgelegter Reihenfolge.
\end{rulebox}

Ein Kampf folgt immer diesem Ablauf:
\begin{enumerate}
	\item Kampfbeginn
	\item Initiative bestimmen
	\item Kampfrunden ausführen
	\item Kampfende
\end{enumerate}

\begin{rulebox}{Initiative}
	Die \term{Initiative}{initiative} bestimmt die Reihenfolge, in der alle
	Kampfteilnehmer handeln.
\end{rulebox}

Zu Beginn eines Kampfes würfelt jeder Beteiligte eine Initiative-Probe.
Die genaue Ausgestaltung der Probe hängt von Talent, Ausrüstung oder Effekten ab.

\subsection*{Varianten für die Initiativereihenfolge}
Der Spielleiter wählt zu Kampfbeginn eine der folgenden Varianten.

\subsubsection*{Variante 1: Klassische Reihenfolge (Einzelinitiative)}
Alle Initiative-Ergebnisse werden von \textbf{hoch nach niedrig} sortiert.
Die Beteiligten handeln in dieser Reihenfolge.

\begin{itemize}
	\item \textbf{Gleichstand:} Zuerst handelt der Charakter mit dem höheren Rang.
	\item Besteht weiterhin Gleichstand, entscheidet der Spielleiter.
\end{itemize}

\subsubsection*{Variante 2: Gruppeninitiative (abwechselnde Hälfte)}
Bei dieser Variante werden die Initiative-Ergebnisse je \textbf{Gruppe} addiert
(z.\,B. „Helden“ gegen „Gegner“).

\begin{enumerate}
	\item Beide Gruppen bilden ihre \textbf{Initiativesumme}.
	\item Die Gruppe mit der höheren Summe ist zuerst am Zug und aktiviert
	\textbf{die Hälfte ihrer Mitglieder} (aufgerundet) in beliebiger Reihenfolge.
	\item Danach aktiviert die andere Gruppe \textbf{die Hälfte ihrer Mitglieder}
	(aufgerundet) in beliebiger Reihenfolge.
	\item Anschließend aktiviert die zuerst gestartete Gruppe ihre \textbf{restlichen}
	Mitglieder, danach die andere Gruppe ihre \textbf{restlichen} Mitglieder.
	\item In den weiteren Runden darf die Reihenfolge der Mitglieder wieder neu gewählt werden.
\end{enumerate}

\begin{rulebox}{Überlegenheit}
	Ist die Initiativesumme einer Gruppe \textbf{mindestens doppelt so hoch} wie die
	der anderen, darf die überlegene Gruppe \textbf{vollständig zuerst} agieren.
\end{rulebox}
 

\section{Kampfrunden}

\begin{rulebox}{Kampfrunde}
	Eine \term{Kampfrunde}{kampfrunde} stellt einen kurzen Zeitraum im Kampf dar, in dem alle
	Beteiligten die Möglichkeit haben zu handeln.
\end{rulebox}

Innerhalb einer Kampfrunde ist jeder Charakter einmal am Zug, wenn er an der
Reihe ist.

\section{Aktionen und Bonusaktionen}

\begin{rulebox}{Aktionen}
	Während seines Zuges kann ein Charakter \textbf{eine \term{Aktion}{aktion}} ausführen.
\end{rulebox}

Typische Aktionen sind:
\begin{itemize}
	\item Angriff mit einer Waffe
	\item Bewegen oder Sprinten
	\item Wirken einer Fähigkeit
	\item Einsatz eines Talents
	\item Eine komplexe Interaktion (z.\,B. Hebel umlegen, Objekt benutzen)
\end{itemize}

\begin{rulebox}{Bonusaktionen}
	Zusätzlich zur Aktion kann ein Charakter pro Runde \textbf{eine \term{Bonusaktion}{bonusaktion}}
	ausführen, sofern er über eine entsprechende Fähigkeit verfügt.
\end{rulebox}

Bonusaktionen sind:
\begin{itemize}
	\item Waffenwechsel
	\item Tränke nutzen
	\item Wirken einer Fähigkeit
\end{itemize}

\section{Bewegung}

\begin{rulebox}{Bewegung im Kampf}
	\term{Bewegung}{bewegung} ist Teil deines Zuges und \textbf{zählt nicht als Aktion}.
\end{rulebox}

Die \textbf{Bewegungsweite} hängt von deinem Talent \textbf{\termref[Gewandtheit]{gewandtheit}} ab.
Die Standard-Bewegungsweite beträgt:
\[
4 + \text{Rang} \text{ Meter}
\]
Gelände und Effekte können diese Distanz erhöhen oder verringern.

\begin{rulebox}{Sprinten im Kampf}
	Mit \term{Sprinten}{sprinten} kannst du als \textbf{Bonusaktion} sprinten und dich
	bis zum \textbf{1,5-fachen} deiner Bewegungsweite bewegen. Zusätzlich musst du eine
	\textbf{Gewandtheitsprobe} bestehen.
\end{rulebox}


\section{Angriff und Verteidigung}

\begin{rulebox}{Grundprinzip}
	Ein \term{Angriff}{angriff} ist eine Aktion, mit der du versuchst, einem Gegner Schaden zuzufügen.
	Der Verteidiger reagiert mit \term{Verteidigung}{verteidigung} durch \textbf{parieren} oder \textbf{ausweichen}.
\end{rulebox}

\subsection*{Würfelgröße}
Die verwendete Würfelgröße ergibt sich aus dem \textbf{Rang der eingesetzten Fähigkeit}.
Bei \textbf{Waffenangriffen oder waffenlosen Angriffen} richtet sich die Würfelgröße nach dem Rang des passenden Talents:
\begin{itemize}
	\item \textbf{Nahkampf:} Talent \textbf{\termref[Kampf]{kampf}}
	\item \textbf{Fernkampf:} Talent \textbf{\termref[Gewandtheit]{gewandtheit}}
\end{itemize}

\subsection*{Ablauf}
Ein Angriff wird immer in diesen Schritten abgehandelt:
\begin{enumerate}
	\item Der Angreifer würfelt den \textbf{Angriffswurf} und bildet seine \textbf{Wurfzahl (Angriff)}.
	\item Der Verteidiger entscheidet sich für eine \textbf{Parade} (Probe mit Talent \termref[Kampf]{kampf}) oder \textbf{Ausweichen} (Probe mit Talent \termref[Gewandtheit]{gewandtheit}), würfelt und bildet seine \textbf{Wurfzahl (Verteidigung)}.
	\item Der Angriff trifft nur, wenn:
	\[
	\text{Wurfzahl (Angriff)} > \text{Wurfzahl (Verteidigung)}
	\]
	\item Trifft der Angriff, wird anschließend der \textbf{\termref[Schaden]{schaden}} ermittelt.
\end{enumerate}

Sowohl Angriff als auch Verteidigung können durch situative Umstände,
Gelände oder aktive Effekte Boni oder Mali erhalten.

\subsection*{Parade vs. Ausweichen}
\term{Parade}{parade} und \term{Ausweichen}{ausweichen} sind die möglichen Reaktionen auf einen Angriff. Eine erfolgreiche Parade ermöglicht einen \textbf{Angriff ohne Probe mit halben \termref[Waffenschaden]{waffenschaden}}. Erfolgreiches Ausweichen ermöglicht eine halbe \termref[Bewegung]{bewegung}.

\subsection*{Magische Waffen bei Paraden}

\begin{rulebox}{Magische Waffen und Parade}
	Magische Waffen haben bei der \textbf{Parade} einen Vorteil gegenüber nicht-magischen Waffen:
	\textbf{+2 Bonus auf die Wurfzahl der Parade}.
\end{rulebox}

Dabei gilt:
\begin{itemize}
	\item Der Angriff oder die Parade muss zuvor \textbf{regulär erfolgreich} sein (Wurfzahl $\ge 4$).
	\item Führen \textbf{beide} Parteien eine magische Waffe, \textbf{hebt sich der Bonus auf} (kein Bonus).
\end{itemize}

\begin{examplebox}{Beispiel: Magische Waffe in der Parade}
	Ein Angreifer führt eine nicht-magische Waffe und würfelt beim Angriff eine 6 (Wurfzahl 6 mit einem W10 - also kein Ass).
	Der Verteidiger pariert mit einer magischen Waffe und erhält +2 auf seine Parade-Wurfzahl.
	Er muss die Wurfzahl des Angreifers erreichen und benötigt daher mindestens eine 4,
	um erfolgreich zu parieren.
\end{examplebox}


\section{Schaden}

\begin{rulebox}{Schaden}
	Erleidet ein Charakter \term{Schaden}{schaden}, werden ihm Punkte von seinen Schutz- und
	Lebensressourcen abgezogen.
\end{rulebox}

Welche Ressource betroffen ist, hängt von Art und Quelle des Schadens ab.
\begin{rulebox}{Reihenfolge der Schadensaufnahme}
	Sofern nicht anders angegeben, wird Schaden in dieser Reihenfolge abgetragen:
	\[
	\textbf{Magischer Schild (SP)} \rightarrow \textbf{Rüstung (RP)} \rightarrow \textbf{Lebenspunkte (LP)}
	\]
\end{rulebox}

\section{Rüstung und Magischer Schild}

\begin{rulebox}{Rüstung (RP)}
	\term{Rüstung}{ruestung} nimmt Schaden auf und reduziert so den Verlust von Lebenspunkten.  
	\textbf{1 Rüstungspunkt (RP)} kann \textbf{2 Schadenspunkte} aufnehmen.
	
	Rüstung kann in der Regel nur \textbf{außerhalb des Kampfes} wiederhergestellt werden.
	Bestimmte Fähigkeiten können \textbf{temporäre} Rüstungspunkte verleihen.
\end{rulebox}

\begin{rulebox}{Magischer Schild (SP)}
	Der \term{Magische Schild}{magischer-schild} nimmt Schaden auf und kann auch während
	des Kampfes erneuert werden.  
	\textbf{1 Schildpunkt (SP)} kann \textbf{1 Schadenspunkt} aufnehmen.
	
	Magischer Schild kann durch Fähigkeiten \textbf{während des Kampfes} wiederhergestellt werden.
\end{rulebox}

% =========================
% Rüstungstypen & Qualitäten
% =========================

\subsection*{(Magische) Rüstungen und magische Roben}

Rüstungen und Roben existieren in unterschiedlichen \textbf{Arten} und \textbf{Qualitäten}.
Sie verleihen sowohl \textbf{Rüstungspunkte (RP)} als auch \textbf{Magischen Schild (SP)} (bei magischen Varianten).

\medskip

\begin{center}
	\begin{tabular}{llcc}
		\toprule
		\textbf{Art} & \textbf{Qualität} & \textbf{RP} & \textbf{SP} \\
		\midrule
		Leichte Rüstung/Robe & schlecht     & 1 & 3 \\
		& normal       & 2 & 4 \\
		& meisterlich  & 3 & 5 \\
		\midrule
		Mittlere Rüstung/Robe & schlecht     & 2 & 4 \\
		& normal       & 3 & 5 \\
		& meisterlich  & 4 & 6 \\
		\midrule
		Schwere Rüstung/Robe & schlecht     & 3 & 5 \\
		& normal       & 4 & 6 \\
		& meisterlich  & 5 & 7 \\
		\bottomrule
	\end{tabular}
\end{center}

\medskip

\textbf{Hinweis zu Rüstungen:}
Die in der Tabelle angegebenen \textbf{SP-Werte} gelten nur für
\textbf{magische Varianten}. Nichtmagische Rüstungen verleihen ausschließlich \textbf{Rüstungspunkte (RP)}.

\medskip
\textbf{Hinweis zu Roben:}
Magische Roben haben nur \textbf{magisches Schild (SP)}, besitzen jedoch häufiger zusätzliche Effekte oder Modifikatoren als magische Rüstungen.

% =========================
% Zusätzliche Rüstungsteile
% =========================

\subsection*{Zusätzliche Rüstungsteile}

Zusätzliche Rüstungsteile (z.\,B. Helm, Armschienen, Beinschienen) verleihen  \textbf{Rüstungspunkte (RP)}.
Nur magische Varianten dieser Teile verleihen zusätzlich \textbf{magisches Schild (SP)}.

\medskip

\begin{center}
	\begin{tabular}{lcc}
		\toprule
		\textbf{Qualität} & \textbf{RP} & \textbf{SP} \\
		\midrule
		schlecht    & 1 & 3 \\
		normal      & 2 & 4 \\
		meisterlich & 3 & 5 \\
		\bottomrule
	\end{tabular}
\end{center}

\medskip

\begin{rulebox}{Kumulative Boni}
	Die Boni mehrerer Rüstungsteile sind \textbf{kumulativ}.
	
	Eine mögliche \textbf{Obergrenze} oder Einschränkung
	liegt im Ermessen der Spielleitung.
\end{rulebox}

% =========================
% Beispiel
% =========================

\begin{examplebox}{Rüstungsberechnung}
	Grundschutz: \textbf{3 RP} \\
	Helm (normal): \textbf{+2 RP}
	
	\medskip
	\textbf{Gesamt: 5 RP}
\end{examplebox}

\section{Schadenstypen}

\term{Schadenstypen}{schadenstypen} bestimmen, wie Schaden verrechnet wird und ob zusätzliche
Statuseffekte ausgelöst werden können. Die Wahrscheinlichkeit, einen Statuseffekt auszulösen, entspricht dem Rang des jeweiligen Effekts.


\begin{center}
	\begin{tabular}{lp{9cm}}
		\toprule
		\textbf{Schadenstyp} & \textbf{Regelwirkung} \\
		\midrule
		\term{Waffenschaden}{waffenschaden} &
		Normaler Schaden durch Waffen oder waffenlose Angriffe
		(SP $\rightarrow$ RP $\rightarrow$ LP) \\
		
		\term{Durchschlagsschaden}{durchschlagsschaden} &
		Ignoriert \termref[Rüstung (RP)]{ruestung}, \termref[Magischen Schild (SP)]{magischer-schild} wirkt normal \\
		
		\term{Direktschaden}{direktschaden} &
		Ignoriert \termref[Rüstung (RP)]{ruestung} und \termref[Magischen Schild (SP)]{magischer-schild} ,
		Schaden wird direkt von den \termref[Lebenspunkten (LP)]{lebenspunkte-lp}  abgezogen \\
		
		\term{Giftschaden}{giftschaden} &
		Kann den Statuseffekt \termref[Vergiftung]{vergiftung}  auslösen \\
		
		\term{Feuerschaden}{feuerschaden} &
		Kann den Statuseffekt \termref[Verbrennung]{verbrennung}  auslösen \\
		
		\term{Blitzschaden}{blitzschaden} &
		Kann den Statuseffekt \termref[Betäubung]{betaeubung}  auslösen \\
		
		\term{Kälteschaden}{kaelteschaden} &
		Kann den Statuseffekt \termref[Unterkühlung]{unterkuehlung}  auslösen \\
		
		\term{Verwesungsschaden}{verwesungsschaden} &
		Kann den Statuseffekt \termref[Erosion]{erosion}  auslösen \\
		\bottomrule
	\end{tabular}
\end{center}

Ob ein Statuseffekt ausgelöst wird, richtet sich nach dem \textbf{Rang}
des verursachenden Schadens (z.B. Verwesungsschaden Rang 3 = 1W8 Chance auf Erosion). Feuer-, Blitz- und Kälteschaden gelten als \term{Elementarschaden}{elementarschaden}.


\section{Statuseffekte}

\begin{rulebox}{Statuseffekte}
	\term{Statuseffekte}{statuseffekte} sind Zustände mit einem eigenen Rang,
	die einen Charakter zeitweise oder dauerhaft beeinflussen.
\end{rulebox}

Wird derselbe Statuseffekt mehrfach angewendet, \textbf{stapelt} er sich:
\begin{itemize}
	\item Jede Anwendung zählt als eigener Effekt.
	\item Effekte werden separat abgehandelt (Schaden, Mali und Dauer gelten jeweils einzeln).
\end{itemize}

\subsection*{Schaden über Zeit}

\begin{center}
	\begin{tabular}{lp{9cm}}
		\toprule
		\textbf{Status} & \textbf{Wirkung} \\
		\midrule
		\term{Vergiftung}{vergiftung} &
		\termref[Direktschaden]{direktschaden} in Höhe des Rangs pro Runde (für die Wirkungsdauer) \\
		
		\term{Verbrennung}{verbrennung} &
		\termref[Schaden]{schaden} in Höhe des Rangs pro Runde (für die Wirkungsdauer) \\
		
		\term{Blutung}{blutung} &
		\termref[Schaden]{schaden} in Höhe von Rang/2 pro Runde; \textbf{steigt jede Runde um +1} (für die Wirkungsdauer) \\
		\bottomrule
	\end{tabular}
\end{center}

\subsection*{Kontroll- und Einschränkungseffekte}

\begin{center}
	\begin{tabular}{lp{9cm}}
		\toprule
		\textbf{Status} & \textbf{Wirkung} \\
		\midrule
		\term{Unterkühlung}{unterkuehlung} &
		Verlust der \termref[Bonusaktion]{bonusaktion} für Rang Runden \\
		
		\term{Betäubung}{betaeubung} &
		Ziel verliert für 1 Runde alle \termref[Aktionen]{aktion} und \termref[Bonusaktionen]{bonusaktion} \\
		
		\term{Blendung}{blendung} &
		Ziel ist 2 Runden geblendet: Malus in Höhe des Rangs auf alle \termref[Aktionen]{aktion} \\
		
		\term{Erschöpfung}{erschoepfung} &
		\termref[Malus]{malus} von -2 für 1 Runde auf Proben mit dem Talent \termref[Gewandtheit]{gewandtheit} \\
		
		\term{Verwirrung}{verwirrung} &
		\termref[Malus]{malus} von -1 für 1 Runde auf Proben mit dem Talent \termref[Kampf]{kampf} \\
		\bottomrule
	\end{tabular}
\end{center}

\subsection*{Permanente Effekte}

\begin{center}
	\begin{tabular}{lp{9cm}}
		\toprule
		\textbf{Status} & \textbf{Wirkung} \\
		\midrule
		\term{Erosion}{erosion} &
		Verlust von Rang $\times$ 1W4 der \textbf{maximalen \termref[Lebenspunkte (LP)]{lebenspunkte-lp}} pro Runde (für die Wirkungsdauer) \\
		
		\term{Heilung}{heilung} &
		Stellt Rang \termref[Lebenspunkte (LP)]{lebenspunkte-lp} wieder her (für die Wirkungsdauer) \\
		
		\term{Entwaffnet}{entwaffnet} &
		Die aktuell geführte Waffe kann nicht eingesetzt werden (bis der Effekt endet) \\
		\bottomrule
	\end{tabular}
\end{center}

\section{Zeitpunkt von Schaden und Statuseffekten}

\begin{rulebox}{Zeitpunkt der Abhandlung}
	\textbf{Angriffsschaden} wird sofort berechnet, sobald ein Angriff trifft und der Schaden ermittelt wird.
	
	\textbf{Schaden und weitere Effekte durch Statuseffekte} werden \textbf{immer zu Beginn einer Runde} abgehandelt.
\end{rulebox}

Dabei gilt:
\begin{itemize}
	\item Effekte, die \textbf{Schaden pro Runde} verursachen (z.\,B. Vergiftung, Verbrennung, Blutung), ticken zu Rundenbeginn.
	\item Effekte, die \textbf{Mali} oder \textbf{Einschränkungen} geben (z.\,B. Blendung, Unterkühlung), gelten ab sofort (noch vor dem nächsten Zug) bis sie auslaufen.
\end{itemize}


\section{Kampfende}

\begin{rulebox}{Ende des Kampfes}
	Ein Kampf endet, wenn keine gegnerischen Handlungen mehr stattfinden oder der
	Konflikt anderweitig aufgelöst ist.
\end{rulebox}

Nach dem Kampf können Effekte auslaufen, Ressourcen regeneriert werden oder
neue Szenen beginnen.
